% Part: lambda-calculus
% Chapter: introduction
% Section: primitive-recursion

\documentclass[../../../include/open-logic-section]{subfiles}

\begin{document}

\olfileid{lam}{int}{pr}
\olsection{\usetoken{S}{lambda definable} فنکشن بنیادی بازگشت ہِٹھ بند نیں}

ہُن بنیادی بازگشت لئی سانوں کجھ کم کرنا پئے گا۔ اسیں ایہ کم
پڑاؤواں وچ کراں گے۔ پہلاں وانگ منّ لئیے کہ ساڈے کول ٹرم $G$ تے $H$
پہلاں ای نیں، جیہڑے فنکشن $g$ تے~$h$ نوں بالترتیب !!{lambda define} کردے نیں۔
سانوں اوہ ٹرم $H$ چاہیدا اے جیہڑا ہِٹھ دِتّے فنکشن~$f$ نوں
!!{lambda define}s:
\begin{align*}
f(0, \vec z) & = g(\vec z) \\
f(x+1, \vec z) & = h(z, f(x,\vec z), \vec z).
\end{align*}
سو عام طور اُتے، جے لامبڈا ٹرم $G'$ تے~$H'$ دِتّے ہون، تاں
ایہو جیہا ٹرم~$F$ لبھنا کافی اے جیہدے لئی
\begin{align*}
F(\num{0}, \vec z) & \equiv G(\vec z) \\
F(\overline{n+1}, \vec z) & \equiv H(\num{n}, F(\num{n}, \vec z), \vec z)
\end{align*}
ہر قدرتی عدد~$n$ لئی ہووے۔ $G'$ تے~$H'$ فنکشن $g$ تے~$h$ نوں
!!{lambda define} کردے نیں، ایس لئی جدوں وی اسیں چرچ عدد $\num{\vec m}$
نوں $\vec z$ دی تھاں رکھاں گے، $F(\num{n+1}, \num{\vec m})$
تقلیل ہو کے صحیح جواب والے نارمل روپ تک پُہنچے گا۔

پر ایس کم لئی ایہو جیہا ٹرم~$F$ لبھنا کافی اے جیہڑا ایہ شرطاں
پوریاں کرے:
\begin{align*}
F(\num 0) & \equiv G \\
F(\overline {n+1}) & \equiv H(\num{n},F(\num{n}))
\intertext{ہر قدرتی عدد $n$ لئی، جتھے}
G & = \lambd[\vec z][G'(\vec z)] \text{ تے}\\
H(u,v) & = \lambd[\vec z][H'(u,v(u,\vec z),\vec z)].
\end{align*}
دوجے لفظاں وچ، لامبڈا دی ایس ترکیب نال وادھو پیرامیٹراں $\vec z$
دی وکھری فکر نہیں کرنی پیندی۔ اوہ لامبڈا لکھت وچ ای سما جاندے نیں۔

ٹرم $F$ دی تعریف توں پہلاں سانوں ترتیب وار جوڑیاں نوں ورتن دا
طریقہ چاہیدا اے۔ اگلا لمہ ایہ طریقہ دیندا اے۔

\begin{lem}
اک لامبڈا ٹرم $D$ موجود اے جیہدے لئی لامبڈا ٹرماں $M$ تے~$N$
دے ہر جوڑے اُتے $D(M,N)(\num{0}) \red M$ تے $D(M,N)(\num{1}) \red N$۔
\end{lem}

\begin{proof}
پہلاں لامبڈا ٹرم $K$ دی تعریف ایویں کرو:
\[
K(y) = \lambd[x][y].
\]
دوجے لفظاں وچ، $K$ اوہ ٹرم $\lambd[y][\lambd[x][y]]$ اے۔ ہور
طرح ویکھیے تاں ہر $M$ لئی، $K(M)$ اک مستقل فنکشن اے جیہڑا
ہر انپٹ اُتے~$M$ واپس کردا اے۔

ہُن $D(x,y,z)$ دی تعریف $D(x,y,z) = z (K(y))x$ نال کرو۔ فیر
ساڈے کول
\begin{align*}
D(M,N,\num 0) & \red \num 0 (K(N)) M \red M \text{ تے}\\
D(M,N,\num 1) & \red \num 1 (K(N)) M \red K(N) M \red N,
\end{align*}
جیویں سانوں چاہیدا سی۔
\end{proof}

خیال ایہ اے کہ $D(M,N)$ جوڑے $\tuple{M, N}$ دی نمائندگی کردا اے۔
جے منّ لئیے کہ $P$ ایہو جیہے جوڑے دی نمائندگی کردا اے، تاں $P(\num 0)$
تے $P(\num 1)$ بالترتیب کھبے تے سجے حصے نوں چُندے نیں، یعنی
$(P)_0$ تے $(P)_1$۔ اسیں ایہ آخری دو لکھتاں ورتاں گے۔

\begin{lem}
!!{lambda definable} فنکشن بنیادی بازگشت ہِٹھ بند نیں۔
\end{lem}

\begin{proof}
سانوں ایہ دکھاؤنا اے کہ کوئی وی ٹرم $G$ تے $H$ دِتّے ہون، اسیں
ایہو جیہا ٹرم $F$ لبھ سکدے آں جیہدے لئی
\begin{align*}
F(\num{0}) & \equiv G \\
F(\num{n+1}) & \equiv H(\num{n}, F(\num{n}))
\end{align*}
ہر قدرتی عدد~$n$ لئی ہووے۔ موٹا خیال ایہ اے کہ چرچ اعداں نوں
وار وار عمل کرن والے ٹرم سمجھ کے \emph{جوڑیاں} دیاں لڑیّاں دا حساب کریے:
\[
\tuple{\num 0, F(\num 0)}, \tuple{\num 1, F(\num 1)}, \dots,
\]
دھیان دیو کہ پہلا جوڑا بس $\tuple{\num 0, G}$ اے۔ جے جوڑا
$\tuple{\num{n}, F(\num{n})}$ دِتّا ہووے، تاں اگلے جوڑے
$\tuple{\num{n+1}, F(\num{n+1})}$ نوں
$\tuple{\num{n+1}, H(\num{n}, F(\num{n}))}$ دے ہم ارز ہونا چاہیدا اے۔
اسیں اک لامبڈا ٹرم~$T$ بناواں گے جیہڑا ایہ اک قدم دا بدل کردا اے۔

تفصیل ایہ اے۔ $T(u)$ دی تعریف ایویں کرو:
\[
T(u) = \tuple{S((u)_0), H((u)_0,(u)_1)}.
\]
ہُن ایہ پرکھنا سوکھا اے کہ کسے وی عدد~$n$ لئی،
\[
T(\tuple{\num{n}, M}) \red \tuple{\num{n+1}, H(\num{n}, M)}.
\]
جیویں اُتے دَسّیا سی، $G$ تے $H$ دِتّے ہون تاں~$F(u)$ دی تعریف
ایویں کرو:
\[
F(u) = (u(T,\tuple{\num 0, G}))_1.
\]
دوجے لفظاں وچ، انپٹ~$\num{n}$ اُتے، $F$، $T$ نوں $n$ واری
$\tuple{\num 0, G}$ توں شروع کر کے لاگو کردا اے، تے فیر دوجا جزو
واپس کردا اے۔ شروع وچ ساڈے کول ایہ اے:
\begin{enumerate}
\item $\num{0} (T, \tuple{\num 0, G}) \equiv \tuple{\num{0}, G}$
\item $F(\num{0}) \equiv G$
\end{enumerate}
~$n$ اُتے استقرا نال اسیں ہر قدرتی عدد لئی ایہ دکھا سکدے آں:
\begin{enumerate}
\item $\num{n+1}(T, \tuple{\num{0}, G}) \equiv \tuple{\num{n+1}, F(\num{n+1})}$
\item $F(\num{n+1}) \equiv H(\num{n}, F(\num{n}))$
\end{enumerate}
دوجی گل لئی ساڈے کول ایہ اے:
\begin{align*}
F(\num{n+1}) & \red (\num{n+1}(T, \tuple{\num 0, G}))_1 \\
& \equiv (T(\num{n} (T, \tuple{\num 0, G})))_1 \\
& \equiv (T(\tuple{\num{n}, F(\num{n})}))_1 \\
& \equiv (\tuple{\num{n+1}, H(\num{n}, F(\num{n}))})_1 \\
& \equiv H(\num{n}, F(\num{n})).
\end{align*}
ایتھے اسیں آخری توں پہلی سطر وچ استقرا دا مفروضہ ورتیا اے۔ پہلی
گل لئی ساڈے کول ایہ اے:
\begin{align*}
\num{n+1} (T, \tuple{\num 0, G}) &
\equiv T(\num{n} (T, \tuple{\num 0, G})) \\
& \equiv T( \tuple{\num{n}, F(\num{n})}) \\
& \equiv \tuple{\num{n+1}, H(\num{n}, F(\num{n}))} \\
& \equiv \tuple{\num{n+1}, F(\num{n+1})}.
\end{align*}
ایتھے اسیں آخری سطر وچ دوجی گل ورتی اے۔ سو اسیں دکھا دِتّا اے کہ
$F(\num 0) \equiv G$ تے ہر $n$ لئی $F(\num {n+1}) \equiv H(\num
n, F(\num{n}))$۔ سانوں عین ایہو چاہیدا سی۔
\end{proof}

\end{document}
